%!TEX program = lualatex
\documentclass[ngerman,headheight=70pt]{scrartcl}
\usepackage[ngerman]{babel}
\usepackage{blindtext}
\usepackage[utf8]{luainputenc}
\usepackage{scrlayer-scrpage}
\usepackage[T1]{fontenc}
\usepackage{fontspec}
\usepackage{multicol}
\usepackage{enumitem}
\usepackage{epigraph}
\usepackage[german=quotes]{csquotes}
\usepackage[bottom=2cm,footskip=8mm]{geometry}
\setmainfont{Source Sans Pro}
\MakeOuterQuote{"}

\setlength{\parindent}{0pt}
\setlength{\parskip}{6pt}

\pagestyle{plain.scrheadings}
\thispagestyle{scrheadings}
\lohead{\vspace*{-1.5cm}\includegraphics[scale=0.4]{../images/stupa_logo_rot.pdf}}
\rohead{\\Präsidium des\\Studierendenparlaments}

\lofoot[Universität Hamburg · Präsidium des Studierendenparlaments\\\hrule Von-Melle-Park 5 · D-20146 Hamburg]{Universität Hamburg · Präsidium des Studierendenparlaments\\\hrule Von-Melle-Park 5 · D-20146 Hamburg}
\cofoot[]{}
\cohead[\thepage]{}

\renewcommand{\thesubsection}{TOP \arabic{subsection}}
\renewcommand{\thesubsubsection}{\arabic{subsubsection}.}
\setcounter{subsection}{-1}

\begin{document}
    UHH · StuPa-Präsidium · Von-Melle-Park 5 · D-20146 Hamburg

    \section*{Beschluss des Studierendenparlamentes vom 16. Juni 2016}
    \subsection*{Demokratische Bildung durch Wissenschaft statt Akkreditierung!}

    \blockquote[Heidelberger Aufruf gegen die Akkreditierung]{Im Jahr 1999 aber
    kam der Bologna-Prozess, und mit ihm kam die Verkündigung, es dürfe keinen
    Studiengang mehr geben, der nicht zuvor die Gnade einer privatwirtschaftlich
    organisierten Akkreditierungsagentur gefunden habe. So wurde es uns eingeredet,
    und so wurde es von den Landesregierungen in allen Bundesländern oktroyiert.
    Seither drehen sich die Universitäten pausenlos im Akkreditierungsrad. Eine
    Evaluation jagt die nächste.}

    Im März diesen Jahres hat das Bundesverfassungsgericht die Unvereinbarkeit
    des aktuellen Akkreditierungs(un)wesens mit der im deutschen Grundgesetz
    entwickelten Wissenschaftsfreiheit (Art. 5) festgestellt. Damit steht ein
    weiteres Instrument aus dem Bologna-Werkzeugkasten heftig in der öffentlichen
    Kritik und kurz davor, auf seinen Platz als Randnotiz in einem Katalog
    (überkommener) hochschulpolitischer Verirrungen verwiesen zu werden.

    In der Folge der Bologna-Reform wurde das bis dato gebräuchliche System von
    Rahmenprüfungsordnungen, die in Aushandlungsprozessen zwischen Wissenschaft
    und Ministerien erarbeitet wurden, durch ein dezentrales Akkreditierungssystem
    ersetzt. Damit sollte dem wettbewerblichen Umbau des Hochschulwesens Vorschub
    geleistet werden.

    Seitdem gibt es einen durch die Bundesländer eingerichteten Akkreditierungsrat,
    dessen Aufgabe die Akkreditierung sog. Akkreditierungsagenturen ist, welche
    wiederrum Studiengänge ("Programmakkreditierung") oder ganze Hochschulen
    ("Systemakkreditierung") akkreditieren. Das klingt umso wahnsinniger, wenn
    man noch berücksichtigt, dass diese Agenturen privatwirtschaftlich organisiert
    sind und zum Teil bspw. durch Arbeitgeberverbände getragen werden. Ihre Aufgabe
    ist dabei die Beurteilung der Qualität von Studiums, wobei in erster Linie
    dessen Form -- es ist in Credit Points ausgemessen, in Modulen portioniert,
    in Bachelor und Master geteilt und mit Prüfungen überfrachtet -- einer
    Betrachtung unterzogen wird.

    Die Kritik des obersten deutschen Gerichts setzt genau hier an: Diese Praxis
    widerspricht offensichtlich dem Grundsatz des Freiheit von Wissenschaft,
    Forschung und Lehre, da Private ohne gesetzlich hinreichend abgesteckte
    Grenzen über Form und Inhalt des Hochschulstudiums entscheiden. Der Gesetzgeber
    hat die Frage, welche Klammer die derzeit über 17.000 Studiengänge im Gebiet
    der Bundesrepublik verbinden soll, aus der Hand gegeben.

    Dadurch verlieren die demokratischen Hochschulgremien und Landesparlamente
    an legitimer Gestaltungshoheit, wobei sowohl zusätzlich enorme Kosten bei
    den Agenturen anfallen, als auch unzählige Unimitarbeiter*innen mit nichts
    anderem beschäftigt sind, als tausende Seiten Papier zu füllen und Begutachter
    zu bezirzen.

    Vor dem Hintergrund der immer breiter vertretenen Kritik an der Bologna-Reform
    und ihren Begleiterscheinungen, so bspw. die Exzellenzinitiative, gerät die
    Vorstellung "unternehmerischer" Hochschulen, verwertbarer Wissenschaft und
    einem arbeitsmarktorientierten Studium mehr und mehr in die Defensive --
    das zeigt sich auch am Akkreditierungswesen. Bereits in den vergangenen Jahren
    wurde das ursprüngliche Modell der Akkreditierungen durch den Protest von
    Seiten der Studierendenschaft, der Lehrenden und der Hochschulleitungen
    (in Form der HRK) stark geschliffen. Jetzt ist die Zeit reif, sich ihm ganz
    zu entledigen.
\end{document}
